\documentclass[11pt]{article}
\usepackage[margin=1in]{geometry}
\usepackage{amsmath,amssymb}
\usepackage[most]{tcolorbox}
\usepackage{xcolor}

\definecolor{boxframe}{RGB}{30,60,110}
\definecolor{boxbg}{RGB}{248,249,252}
\definecolor{headbg}{RGB}{30,60,110}

\newtcolorbox{geomproblem}[1][]{
    enhanced,
    colback=boxbg,
    colframe=boxframe,
    colbacktitle=headbg,
    coltitle=white,
    fonttitle=\bfseries\large,
    title=Geometric Problem Statement,
    boxrule=1pt,
    arc=2mm,
    top=3mm, bottom=3mm, left=4mm, right=4mm,
    #1
}

\begin{document}
\thispagestyle{empty}

\begin{geomproblem}

\textbf{Given Conditions / Setup}
\begin{itemize}
    \item Let $ABC$ be a triangle with incircle of radius $r$.
    \item Let $P$ be a point on side $BC$ such that $BP = m = 4$ and $PC = n = 3$.
    \item Let $AP = u = 5$ be the cevian from $A$ to $P$.
    \item The cevian $AP$ splits $\triangle ABC$ into $\triangle ABP$ and $\triangle ACP$, with inradii $r_2 = 0.9$ and $r_3 = 0.8$ respectively.
\end{itemize}

\vspace{2mm}
\textbf{Definitions of Parameters}
\begin{align*}
    r &= \text{inradius of the full triangle ABC} \\
    r_2 &= \text{inradius of sub-triangle ABP} \\
    r_3 &= \text{inradius of sub-triangle ACP} \\
    u &= \text{length of cevian AP} \\
    m &= \text{length of segment BP} \\
    n &= \text{length of segment PC}
\end{align*}

\vspace{2mm}
\textbf{Proof Goal}
\begin{center}
\fbox{$\displaystyle \frac{1}{r_2} + \frac{1}{r_3} \;=\; \frac{2}{r} + \frac{u(m+n)}{m\,n\,r}$}
\end{center}

\vspace{3mm}
\hfill\textit{\small Date: September 22, 2026 \quad | \quad Author: Hasin Khan}

\end{geomproblem}

\end{document}
